\documentclass[11pt]{article}
\usepackage[margin=1in]{geometry}
\usepackage{amsmath,amssymb,amsthm}
\usepackage[T1]{fontenc}
\usepackage{lmodern}
\usepackage{graphicx}
\usepackage[hidelinks]{hyperref}
\usepackage{xcolor}

\newtheorem{theorem}{Theorem}
\newtheorem{proposition}[theorem]{Proposition}
\newtheorem{remark}[theorem]{Remark}
\newcommand{\R}{\mathbb R}
\newcommand{\Lp}{L^p}
\newcommand{\MR}{\mathrm{MR}^p}
\newcommand{\Lin}{\mathcal L}

\title{Uniform Propagation Removes the Extra Diagonal Factor\\
in Nonautonomous Maximal Regularity}
\author{Candidate partial result for arXiv:2403.00386}
\date{11 August 2026}

\begin{document}
\maketitle

\begin{center}
\fcolorbox{black}{gray!8}{\parbox{0.88\linewidth}{
\textbf{Status: candidate partial result, likely valid.}
The general question in the source remains open.  The extra polynomial factor
is removed under an additional, natural propagation hypothesis, and the
remaining exponential rate is shown to be optimal.}}
\end{center}

\section{Source question}

Let $I=(a,b)$, $T=b-a<\infty$, let $1<p<\infty$, and let
$D\hookrightarrow X$ be the common domain pair used in the source paper.
For a nonautonomous maximally regular operator $A$, write
\[
 K=[A]_{\mathrm{MR}_p(I)}
   =\|L_A^{-1}\|_{\Lin(\Lp(I;X),\MR_0(I))},
 \qquad L_Au=u'+A(\cdot)u.
\]
Here $\MR(I)=W^{1,p}(I;X)\cap\Lp(I;D)$ has the source paper's
time-homogeneous norm
\[
 \|u\|_{\MR(I)}^p
 =\|u\|_{\Lp(I;X)}^p+\|u\|_{\Lp(I;D)}^p
   +T^p\|u'\|_{\Lp(I;X)}^p,
\]
and the subscript zero denotes $u(a)=0$.

Belin and Lafitte prove
\[
 [A+\lambda]_{\mathrm{MR}_p(I)}
 \le c_p(\lambda T)[A]_{\mathrm{MR}_p(I)},
\]
where $c_p(\nu)\sim-\nu e^{-\nu}$ as $\nu\to-\infty$ and
$c_p(\nu)\sim\nu$ as $\nu\to+\infty$.  Immediately after Theorem 2.5
they ask whether the extra factor $\nu$ can be discarded; Remark 3.4 repeats
that sharpness at either infinity is unknown \cite{BL}.

\begin{figure}[ht]
  \centering
  \includegraphics[width=\textwidth]{figures/open_problem_crop.png}
  \caption{Theorem 2.5 and the open sharpness remark in
  arXiv:2403.00386v2, PDF page 14.}
\end{figure}
\clearpage

\section{Proof intuition}

The source proof uses the scalar conjugation
$u(t)=e^{-\lambda(t-a)}v(t)$ and estimates the full maximal-regularity norm
of $v$ before undoing the conjugation.  Differentiating the scalar weight
then produces the extra factor $|\lambda|T$.

If the unshifted problem has a uniformly $X$-bounded propagator, one can
instead estimate only the $X$-norm of the shifted solution.  Its Green kernel
is the old kernel multiplied by $e^{-\lambda(t-s)}$, so Young's inequality
gives the exact scalar $L^1$ factor.  Returning to the unshifted equation,
\[
 L_Au=f-\lambda u,
\]
then controls the entire maximal-regularity norm at once.  This avoids
differentiating the exponential weight inside the norm and hence avoids the
extra polynomial factor.

\section{Factor removal under uniformly bounded propagation on X}

We isolate precisely the additional hypothesis used below.

\begin{theorem}[Factor-free diagonal estimate]\label{thm:main}
Assume that $A\in\mathrm{MR}_p(I)$ and that its zero-initial solution
operator $R_0=L_A^{-1}$ has a Green representation
\[
 (R_0f)(t)=\int_a^t U(t,s)f(s)\,ds
 \quad\text{for almost every }t\in I,
 \tag{1}
\]
where $U(t,s)\in\Lin(X)$ is strongly measurable on $a<s<t<b$ and
\[
 M:=\operatorname*{ess\,sup}_{a<s<t<b}\|U(t,s)\|_{\Lin(X)}<\infty.
 \tag{2}
\]
Then, for every real $\lambda$,
\[
 \boxed{
 [A+\lambda]_{\mathrm{MR}_p(I)}
 \le \bigl(1+M|1-e^{-\lambda T}|\bigr)
       [A]_{\mathrm{MR}_p(I)}.}
 \tag{3}
\]
Consequently
\[
 [A+\lambda]_{\mathrm{MR}_p(I)}=O(1)
       \quad(\lambda\to+\infty),
 \qquad
 [A-\mu]_{\mathrm{MR}_p(I)}=O(e^{\mu T})
       \quad(\mu\to+\infty).
\]
Thus the extra factor $|\lambda|T$ disappears in both regimes.
\end{theorem}

\begin{proof}
Fix $f\in\Lp(I;X)$.  The usual scalar conjugation, which is valid on the
finite interval, shows that the unique solution of
\[
 u'+(A(t)+\lambda)u=f,\qquad u(a)=0,
 \tag{4}
\]
is
\[
 u(t)=e^{-\lambda(t-a)}
 R_0\bigl(e^{\lambda(\,\cdot-a)}f\bigr)(t).
\]
Using (1), this becomes
\[
 u(t)=\int_a^t e^{-\lambda(t-s)}U(t,s)f(s)\,ds.
 \tag{5}
\]
Extend $f$ by zero outside $I$ and define
\[
 k_\lambda(r)=e^{-\lambda r}\mathbf 1_{(0,T)}(r).
\]
Equations (2)--(5), followed by the scalar Young convolution inequality,
give
\[
 \|u\|_{\Lp(I;X)}
 \le M\|k_\lambda\|_{L^1(\R)}\|f\|_{\Lp(I;X)}
 =M h_\lambda(T)\|f\|_{\Lp(I;X)},
 \tag{6}
\]
where
\[
 h_\lambda(T)=
 \begin{cases}
 (1-e^{-\lambda T})/\lambda,&\lambda\ne0,\\
 T,&\lambda=0.
 \end{cases}
\]
This quantity is positive for every real $\lambda$.

Now regard (4) as the unshifted equation
$L_Au=f-\lambda u$.  By the definition of
$K=[A]_{\mathrm{MR}_p(I)}$ and (6),
\begin{align*}
 \|u\|_{\MR(I)}
 &\le K\bigl(\|f\|_{\Lp(I;X)}
       +|\lambda|\|u\|_{\Lp(I;X)}\bigr)\\
 &\le K\bigl(1+M|\lambda|h_\lambda(T)\bigr)
       \|f\|_{\Lp(I;X)}\\
 &=K\bigl(1+M|1-e^{-\lambda T}|\bigr)
       \|f\|_{\Lp(I;X)}.
\end{align*}
The same formula at $\lambda=0$ follows directly, and taking the optimal
constant proves (3).  Its two stated asymptotics are immediate.
\end{proof}

\begin{remark}[What the hypothesis means]
Hypothesis (1) is standard when $A$ generates an evolution family on $X$
which represents inhomogeneous zero-initial solutions, and (2) is uniform
boundedness of that evolution family on the finite time triangle.  The
autonomous estimate in Lemma 3.3 of the source is the special case
$U(t,s)=e^{-(t-s)A}$.  Theorem~\ref{thm:main} therefore extends that exact
factor-free mechanism to a genuinely nonautonomous class.
\end{remark}

\section{The exponential rate cannot be improved}

The next elementary computation shows that the remaining negative-shift
exponential in Theorem~\ref{thm:main} has the correct order even in one
dimension.

\begin{proposition}[Scalar lower bound]\label{prop:sharp}
Let $I=(0,T)$, $X=D=\R$ with the same norm, and $A(t)=0$.  There is a
constant $c_p>0$, depending only on $p$, such that for every $\mu T\ge2$,
\[
 [-\mu]_{\mathrm{MR}_p(0,T)}\ge c_p T e^{\mu T}.
 \tag{7}
\]
In particular no estimate uniform in this class can replace
$e^{\mu T}$ by a smaller exponential order as $\mu T\to\infty$.
\end{proposition}

\begin{proof}
Set
\[
 f_\mu(s)=\mu^{1/p}\mathbf 1_{[0,1/\mu]}(s).
\]
Then $\|f_\mu\|_{L^p(0,T)}=1$.  The solution of
$u'-\mu u=f_\mu$, $u(0)=0$, is
\[
 u(t)=\int_0^t e^{\mu(t-s)}f_\mu(s)\,ds.
\]
For $t\ge1/\mu$ this gives
\[
 u(t)=(1-e^{-1})\mu^{1/p-1}e^{\mu t},
 \qquad
 u'(t)=(1-e^{-1})\mu^{1/p}e^{\mu t}.
\]
Because $[T-1/\mu,T]\subset[1/\mu,T]$ when $\mu T\ge2$,
\begin{align*}
 \|u'\|_{L^p(0,T)}^p
 &\ge (1-e^{-1})^p\mu
       \int_{T-1/\mu}^{T}e^{p\mu t}\,dt\\
 &= (1-e^{-1})^p\frac{1-e^{-p}}p e^{p\mu T}.
\end{align*}
The derivative term in the source norm therefore yields (7) with
\[
 c_p=(1-e^{-1})\left(\frac{1-e^{-p}}p\right)^{1/p}>0.
\]
\end{proof}

\section{Scope and obstruction to the full question}

Theorem~\ref{thm:main} does \emph{not} settle the source question for every
$A\in\mathrm{MR}_p(I)$.  Bare maximal regularity gives the trace embedding
\[
 \MR(I)\hookrightarrow C(\overline I;(D,X)_{1-1/p,p}),
\]
so homogeneous propagation is naturally controlled on the trace space.
It does not by itself supply a Green family uniformly bounded on the rough
space $X$, which is exactly the input needed in (6).  Attempting to use only
the bounded Volterra inverse $R_0$ loses the pointwise time-triangle estimate;
Neumann-series bounds then depend on the maximal-regularity constant in the
exponent and do not recover the scale-sharp factor $e^{|\lambda|T}$.

Thus the remaining full problem can be stated sharply: determine whether
the $X$-space convolution estimate behind (6), or an adequate substitute,
follows from nonautonomous maximal regularity alone.  Alternatively, a
counterexample would need maximally regular coefficients whose propagation
is controlled only on the trace space and for which the additional
$|\lambda|T$ loss is genuinely attained.

The successful theorem was checked by two independent derivations: the
scalar conjugation identity gives (5), while substitution into
$L_Au=f-\lambda u$ gives the maximal-regularity estimate.  Proposition
\ref{prop:sharp} was checked by direct integration.  A bounded search by
arXiv id, exact open-question wording, and the phrases \emph{nonautonomous
maximal regularity}, \emph{diagonal perturbation}, \emph{evolution family},
and \emph{sharp constant} located the source and later versions of the same
paper, but no later resolution of the question or exact statement of
Theorem~\ref{thm:main}.  This is evidence, not a guarantee, of novelty.

\begin{thebibliography}{9}
\bibitem{BL}
T. Belin and P. Lafitte,
\emph{Quantitative estimates of $L^p$ maximal regularity for nonautonomous
operators and global existence for quasilinear equations},
arXiv:2403.00386v2 (2024); Electron. J. Differential Equations 2025, no.~18.
\end{thebibliography}

\end{document}
